% page 80 of the facsimile (image p-079 of the scan)
% block 160.0 x 233.3 mm ; left margin 8.0 mm
\pgc{-12.700mm}{0.000mm}{
\l{\cel{6}72}
\l{}
\l{}
\l{\sou{brakiopodo}. (\sou{zool.}) Klaso de moluski kun konko bivalva,}
\l{\cel{3}sen-movila, di qua la boko situesas inter du brakii tre}
\l{\cel{3}karnoza ed esas provizita ye multa filamenti quin oli}
\l{\cel{3}povas igar ekirar od enirar per spular olci spirale. DEFIS.}
\l{}
\l{\sou{brakistokrono}.\cel{2}(\sou{Mekan.}) Kurvo quan devas sequar korpo pezoza}
\l{\cel{3}por irar de ca a ta punto en tempo minima, la rapideso-grado}
\l{\cel{3}lor la departo esante nula. - DEFI.}
\l{}
\l{\sou{brakteo}. (\sou{bot.}) Folieto qua diferas de la folio ordinara qua,}
\l{\cel{3}che ula planti, situite ye la inserto-punto di la floro,}
\l{\cel{3}akompanas o parte kovras olca. - DEFIS. - DEFIS.}
\l{}
\l{\sou{bramar}.\cel{2}(\sou{netrans}.) Vorto generala por omna krii di bestii.}
\l{\cel{3}- FIS.}
\l{}
\l{\sou{bramano}. Sacerdoto di la bramanismo, qua apartenas a la unesma}
\l{\cel{3}ek la kasti Indiana, e qua dicis naskir de la kapo di}
\l{\cel{3}Brama. - DEFIRS.}
\l{}
\l{\sou{brano}. Reziduo de la muelo, asemblajo ek la eskombri de la}
\l{\cel{3}shelo di la frumento-grani. - EF.}
\l{}
\l{\sou{brancho}. I. (\sou{bot}.) Sproso ligna qua kreskas sur la trunko di}
\l{\cel{3}arboro. - II. (\sou{geom}.)Singla de la linei intersimetra,di}
\l{\cel{3}qui la uniono formacas ula kurvi geometriala (hiperbolo,}
\l{\cel{3}parabolo). - III. To omna quo, departante de stango, aspek-}
\l{\cel{3}tas analoge kam la arboro-branchi. - IV. La familii diversa,}
\l{\cel{3}genitita de trunko komuna. (Ex.: familio-branchi). - DEFIR.}
\l{}
\l{\sou{brando}. Fasko de palio flamoza per qua onu dis-tushas por in-}
\l{\cel{3}cendiar, o quan onu kunportas por lumizar. - DEFI.}
\l{}
\l{\sou{brandio}. Kompozajo ek aquo ed alkoholo, obtenita per la distilo}
\l{\cel{3}di vino, cidro, grani, betravo, e c. - DEF.}
\l{}
\l{\sou{brandisar}. (\sou{trans.)}\cel{3}Ociligar en sua manuo minacoze (espado,}
\l{\cel{3}lanco, javelino). - EFIS.}
\l{}
\l{\sou{brankardo.}\cel{2}Singla de la du branchi (ligna), plu o min lejera,}
\l{\cel{3}maxim ofte kurvigita por su adaptar a la korpo di la kavalo,}
\l{\cel{3}qui artike juntesas a la avana parto di veturo, ed inter qui}
\l{\cel{3}onu jungas la kavalo. - F.}
\l{}
\l{\sou{brankio}. Lami dispozita quale la denti di pektilo, ye qua kon-}
\l{\cel{3}sistas la respir-aparato di la animali qui vivas en aquo e}
\l{\cel{3}respiras la aero quan olca kontenas dissolvite (fishi, krus-}
\l{\cel{3}tacei). - EFIS.}
\l{}
\l{\sou{brankiopodo}. (\sou{zool.}) Mikra krustaceo, di qua la pedi, apta por}
\l{\cel{3}natado, esas garnisita ye apendici brankia. - DEFIS.}
}
